\documentclass[12pt,a4paper]{report}

<<<<<<< Updated upstream
% ============================================================
% PREAMBLE  –  compile with: xelatex thesis.tex
% ============================================================

\usepackage{fontspec}         % Native font handling for XeLaTeX
=======
% ================================================================
% PACKAGES
% ================================================================
\usepackage{fontspec}
>>>>>>> Stashed changes
\usepackage{geometry}
\usepackage{graphicx}
\usepackage{fancyhdr}
\usepackage{setspace}
\usepackage{tocloft}
\usepackage{titlesec}
\usepackage{amsmath}
\usepackage{amssymb}
\usepackage{array}
<<<<<<< Updated upstream
\usepackage{booktabs}         % Professional horizontal table rules
%\usepackage{parskip}          % Space between paragraphs instead of indent
\usepackage{microtype}        % Micro-typographic improvements

% ----- THE TWO PACKAGES THAT WERE MISSING AND CAUSED ALL FLOAT ISSUES -----
\usepackage{float}            % REQUIRED for [H] – forces figure/table exactly here
\usepackage{placeins}         % REQUIRED for \FloatBarrier – hard section boundary
% --------------------------------------------------------------------------

% Arabic support (polyglossia replaces babel; the two are incompatible)
=======
\usepackage{booktabs}
\usepackage{longtable}
\usepackage{tabularx}
\usepackage{caption}
\usepackage[ruled,vlined,linesnumbered]{algorithm2e}
>>>>>>> Stashed changes
\usepackage{polyglossia}

% Arabic support for cover page
\setmainlanguage{english}
\setotherlanguage{arabic}
\newfontfamily\arabicfont[Script=Arabic]{Amiri}

% hyperref MUST be loaded last
\usepackage{hyperref}

<<<<<<< Updated upstream
% ------------------------------------------------------------
% Page geometry
% ------------------------------------------------------------
\geometry{
    left=3cm,
    right=2.5cm,
    top=2.5cm,
    bottom=2.5cm
}

% ------------------------------------------------------------
% Line spacing
% ------------------------------------------------------------
\onehalfspacing

% ------------------------------------------------------------
% Header and footer
% ------------------------------------------------------------
=======
% ================================================================
% GEOMETRY & STYLE
% ================================================================
\geometry{left=3cm, right=2.5cm, top=2.5cm, bottom=2.5cm}
\onehalfspacing

>>>>>>> Stashed changes
\pagestyle{fancy}
\fancyhf{}
\fancyhead[L]{\leftmark}
\fancyhead[R]{\thepage}
\renewcommand{\headrulewidth}{0.4pt}

<<<<<<< Updated upstream
% ------------------------------------------------------------
% Chapter and section formatting
% ------------------------------------------------------------
=======
>>>>>>> Stashed changes
\titleformat{\chapter}[display]
  {\normalfont\huge\bfseries}{\chaptertitlename\ \thechapter}{20pt}{\Huge}
\titlespacing*{\chapter}{0pt}{-20pt}{40pt}

<<<<<<< Updated upstream
% ------------------------------------------------------------
% Hyperref setup
% ------------------------------------------------------------
=======
>>>>>>> Stashed changes
\hypersetup{
    colorlinks=true,
    linkcolor=black,
    citecolor=blue,
    urlcolor=blue,
}

<<<<<<< Updated upstream
% ============================================================
=======
% ================================================================
% DOCUMENT START
% ================================================================
>>>>>>> Stashed changes
\begin{document}
% ============================================================

<<<<<<< Updated upstream
% ============================================================
% TITLE PAGE
% ============================================================
\newgeometry{left=2.5cm, right=2.5cm, top=1.5cm, bottom=1.5cm} % Temporary smaller margins for title page
=======
% ----------------------------------------------------------------
% COVER PAGE
% ----------------------------------------------------------------
\newgeometry{left=2.5cm, right=2.5cm, top=1.5cm, bottom=1.5cm}
>>>>>>> Stashed changes
\begin{titlepage}
\thispagestyle{empty}
\centering

\vspace*{-1cm}

{\footnotesize Democratic and People's Republic of Algeria}\\
{\footnotesize \textarabic{الجمهورية الجزائرية الديمقراطية الشعبية}}

{\footnotesize Ministry of Higher Education and Scientific Research}\\
{\footnotesize \textarabic{وزارة التعليم العالي و البحث العلمي}}

\vspace{0.2cm}
\rule{\textwidth}{0.3pt}
\vspace{0.2cm}

\begin{minipage}{0.25\textwidth}
\centering

\includegraphics[width=0.9\textwidth]{usthb.png}
\end{minipage}
\hfill
\begin{minipage}{0.7\textwidth}
\centering
{\footnotesize \textarabic{جامعة العلوم والتكنولوجيا هواري بومدين}}\\
{\footnotesize University of Science and Technology Houari Boumediene}
\end{minipage}

\vspace{0.4cm}

{\LARGE \textbf{Faculty of Informatics}}\\
\vspace{0.5cm}
{\normalsize Thesis for graduation}\\
{\normalsize for the obtainment of the Bachelor degree}\\
\vspace{0.5cm}
{\normalsize \textbf{Option : ACAD}}

\vspace{0.3cm}
\rule{\textwidth}{0.6pt}
\vspace{0.3cm}

{\Large \textbf{Design and Implementation of an Automatic and Intelligent IT Resource Provisioning Solution (Servers and Virtual Machines) for the Scalable Deployment of a High-Load Web Application.}}

\vspace{0.3cm}
\rule{\textwidth}{0.6pt}
\vspace{0.3cm}

\begin{minipage}{0.45\textwidth}
\normalsize \textit{Realized by :}\\
Mr. Zakaria \textsc{HAMMAL}\\
Mr. Oussama \textsc{SELATENIA}\\
Mr. Oussama \textsc{MENNOUR}\\
Mr. Younes \textsc{DJEHA}
\end{minipage}
\hfill
\begin{minipage}{0.45\textwidth}
\normalsize \textit{Supervised by :}\\
Mr. Khaled \textsc{Zeraoulia}
\end{minipage}

\vfill

{\normalsize Held on [DATE], before the jury composed of :}

{\normalsize Mr.\textsc{LASTNAME} Firstname \hspace{0.5cm} -- President of jury}\\
{\normalsize Mr.\textsc{LASTNAME} Firstname \hspace{0.5cm} -- Examiner}

\vspace{0.8cm}

{\normalsize Promotion : 2025/2026}\\
{\normalsize Code PFE : [CODE]}

\end{titlepage}
\restoregeometry

<<<<<<< Updated upstream
% ============================================================
% FRONT MATTER  –  Roman numeral page numbering
% ============================================================
\pagenumbering{roman}

\newpage
=======
% ----------------------------------------------------------------
% FRONT MATTER
% ----------------------------------------------------------------
\pagenumbering{roman}

>>>>>>> Stashed changes
\chapter*{Dedication}
\addcontentsline{toc}{chapter}{Dedication}
\begin{center}
\itshape
[Content derived from the original memoire.tex dedication section.]
\end{center}

<<<<<<< Updated upstream
\newpage
\chapter*{Acknowledgements}
\addcontentsline{toc}{chapter}{Acknowledgements}
To be implemented

\newpage
\chapter*{Abstract}
\addcontentsline{toc}{chapter}{Abstract}
To be implemented
=======
\chapter*{Acknowledgements}
\addcontentsline{toc}{chapter}{Acknowledgements}
[Content derived from the original memoire.tex acknowledgements section.]

\chapter*{Abstract}
\addcontentsline{toc}{chapter}{Abstract}

\section*{English}
The rapid expansion of digital healthcare services necessitates infrastructure that reconciles high availability with operational cost-efficiency. Traditional manual resource allocation fails to respond to the dynamic volatility of high-load web applications. This thesis proposes an autonomous hierarchical framework for intelligent resource provisioning. The architecture integrates a hybrid meta-heuristic algorithm, combining Whale Optimisation (WOA) and Differential Evolution (DE), to optimize virtual machine placement across physical nodes. Furthermore, a dual-agent system, comprising a Planner and a Critic, manages remediation actions based on risk-stratified escalation protocols. The solution is validated through the deployment of a distributed medical information system, demonstrating proactive scalability and resilient fault domain isolation.
>>>>>>> Stashed changes

\section*{Français}
Conception et réalisation d'une solution de provisioning automatique et intelligent de ressources IT (serveurs et machines gratuites) pour le déploiement scalable d'une application web à forte charge. L'expansion des services de santé numériques nécessite des infrastructures conciliant haute disponibilité et efficacité des coûts. Cette thèse propose un cadre hiérarchique autonome pour le provisioning intelligent des ressources. L'architecture intègre un algorithme méta-heuristique hybride (WOA-DE) pour optimiser le placement des machines virtuelles. Un système à double agent (Planner/Critic) gère les actions de remédiation selon des protocoles d'escalade.

<<<<<<< Updated upstream
\newpage
\tableofcontents

\newpage
\listoffigures
\addcontentsline{toc}{chapter}{List of Figures}

\newpage
\listoftables
\addcontentsline{toc}{chapter}{List of Tables}

\newpage
\chapter*{List of Abbreviations}
\addcontentsline{toc}{chapter}{List of Abbreviations}
\begin{tabbing}
\hspace{3cm} \= \kill
\textbf{API}   \> Application Programming Interface \\
\textbf{CI/CD} \> Continuous Integration/Continuous Deployment \\
\textbf{CPU}   \> Central Processing Unit \\
\textbf{DE}    \> Differential Evolution \\
\textbf{DevOps}\> Development and Operations \\
\textbf{HTTP}  \> Hypertext Transfer Protocol \\
\textbf{IT}    \> Information Technology \\
\textbf{OS}    \> Operating System \\
\textbf{REST}  \> Representational State Transfer \\
\textbf{SaaS}  \> Software as a Service \\
\textbf{VM}    \> Virtual Machine \\
\textbf{VPN}   \> Virtual Private Network \\
\textbf{WOA}   \> Whale Optimisation Algorithm \\
\end{tabbing}

\newpage
\chapter*{General Introduction}
\addcontentsline{toc}{chapter}{General Introduction}
To be implemented
\newpage

% ============================================================
% MAIN MATTER  –  Arabic numeral page numbering from 1
% ============================================================
\pagenumbering{arabic}
\setcounter{page}{1}

% ============================================================
% CHAPTER 1 – Theoretical Foundation and Architectural Paradigms
% ============================================================
\chapter{Theoretical Foundation and Architectural Paradigms}
\label{chap:theory}

\section{Introduction}

Modern software development has been shaped by evolving scalability demands and the
need for efficient resource utilization. Virtualization and containerization have
fundamentally altered how applications are developed, deployed, and managed. This
chapter establishes the theoretical foundation for these paradigm shifts, covering
their core concepts, operational characteristics, and relationship to microservices
architecture, before critically examining whether organizations face a genuine
imperative to adopt them.
=======
\tableofcontents
\listoffigures
\listoftables

% ----------------------------------------------------------------
% MAIN MATTER
% ----------------------------------------------------------------
\pagenumbering{arabic}
\setcounter{page}{1}

\chapter*{General Introduction}
\addcontentsline{toc}{chapter}{General Introduction}

The global transition toward digitized healthcare systems has imposed unprecedented demands on the resilience and scalability of IT infrastructures. Modern medical information systems must sustain continuous availability while processing high volumes of sensitive data under volatile traffic patterns. Traditional infrastructure management, characterized by static provisioning and manual oversight, proves insufficient for these high-load environments. The latency associated with human intervention and the inefficiencies of rigid resource allocation often lead to service degradation or excessive capital expenditure.

Achieving a balance between hardware efficiency and service continuity presents a structural tension. Dense workload placement maximizes utilization but increases the risk associated with physical node failures. Conversely, over-provisioning ensures a safety margin but results in significant underutilization of hardware assets. This thesis addresses this problem by proposing an autonomous provisioning framework that utilizes intelligent scheduling and predictive remediation.

The objective of this research is the design and implementation of a self-regulating cluster architecture capable of managing microservices at scale. By leveraging meta-heuristic optimization for virtual machine placement and risk-stratified cognitive agents for orchestration, the system provides a robust solution for deploying complex web applications. The study is organized into three distinct phases: a review of current abstraction technologies, the theoretical architectural design, and the concrete implementation of a distributed medical information system.

\chapter{State of the Art}

\section{Infrastructural Abstraction Paradigms}

Modern software deployment relies on the abstraction of physical hardware into logical units to facilitate multi-tenancy and efficient resource sharing.

\subsection{Hypervisor Classification}
Virtualization allows multiple guest operating systems to execute concurrently on a single physical host. Table \ref{tab:hypervisors} compares the two fundamental hypervisor architectures.

\begin{table}[htbp]
\centering
\caption{Comparison of Type 1 and Type 2 Hypervisors}
\label{tab:hypervisors}
\begin{tabular}{@{}lll@{}}
\toprule
\textbf{Feature} & \textbf{Type 1 (Bare-Metal)} & \textbf{Type 2 (Hosted)} \\ \midrule
Execution & Direct on Hardware & Atop Host OS \\
Performance & High (Near-native) & Moderate (Overhead) \\
Isolation & Strong hardware-level & Dependent on Host OS security \\
Examples & Proxmox VE (KVM), ESXi & VirtualBox, VMware Workstation \\ \bottomrule
\end{tabular}
\end{table}

\subsection{Virtual Machines versus Containerization}
While virtual machines provide comprehensive hardware isolation, containerization virtualizes at the operating system level. Table \ref{tab:vm_vs_container} summarizes their operational characteristics.

\begin{table}[htbp]
\centering
\caption{Architecture: Virtual Machines vs. Containers}
\label{tab:vm_vs_container}
\begin{tabular}{@{}lll@{}}
\toprule
\textbf{Metric} & \textbf{Virtual Machines} & \textbf{Containers} \\ \midrule
Isolation & Hardware-level (Strong) & OS-level (Lightweight) \\
Guest OS & Full OS required & Shared Host Kernel \\
Startup Time & Minutes & Milliseconds \\
Density & Moderate & High \\ \bottomrule
\end{tabular}
\end{table}

\section{Software Architecture Patterns}

The structural design of applications dictates their scalability. Monolithic architectures integrate all components into a single unit, whereas microservices decompose logic into autonomous entities.

\begin{table}[htbp]
\centering
\caption{Monolithic vs. Microservices Comparison}
\label{tab:monolith_vs_micro}
\begin{tabular}{@{}lll@{}}
\toprule
\textbf{Characteristic} & \textbf{Monolithic} & \textbf{Microservices} \\ \midrule
Deployment & Single artifact & Independent artifacts \\
Scaling & Vertical (All components) & Horizontal (Specific pods) \\
Fault Isolation & Low (Cascading failure) & High (Isolated failure) \\
Data Storage & Centralized & Decentralized (DB per service) \\ \bottomrule
\end{tabular}
\end{table}

\section{Related Works}

The problem of optimal workload placement and autonomous provisioning has been addressed using various algorithmic approaches. Table \ref{tab:related_works} summarizes the state of the art.

\begin{longtable}{@{}p{3cm}p{1cm}p{3cm}p{3cm}p{3cm}@{}}
\caption{Summary of Related Works in Resource Management} \label{tab:related_works} \\
\toprule
\textbf{Authors} & \textbf{Year} & \textbf{Methodology} & \textbf{Strengths} & \textbf{Limitations} \\ \midrule
\endfirsthead
\toprule
\textbf{Authors} & \textbf{Year} & \textbf{Methodology} & \textbf{Strengths} & \textbf{Limitations} \\ \midrule
\endhead
\bottomrule
\endfoot
Mirjalili et al. & 2016 & Whale Optimisation (WOA) & Effective exploitation & Premature convergence \\ \midrule
Storn \& Price & 1997 & Differential Evolution (DE) & High diversity & Slow in local search \\ \midrule
Saxena et al. & 2022 & Proactive Scaling & Predictive accuracy & Data intensity \\ \bottomrule
\end{longtable}

\chapter{Conception}

\section{Introduction}

This chapter delineates the theoretical and architectural design of the autonomous cluster management system. The design is technology-agnostic, presenting a generalized model for resilient infrastructure that satisfies the requirements of high-load web applications. It defines the hierarchical layers, the mathematical optimization heuristics, and the cognitive agent logic required for self-regulation.

\section{Design Objectives and Validation Criteria}

The primary objective is the resolution of the structural tension between resource efficiency and service availability. The system must maintain high node density while ensuring performance isolation and headroom for traffic surges. The proposed architecture is validated through four criteria: automated provisioning latency, responsive remediation of resource contention, constraint-respecting scheduling, and the preservation of continuous availability during failure events.

\section{Layered Infrastructure Architecture}

The infrastructure is conceptualized as a four-layer hierarchical stack. Each layer provides abstracted services to the layer above. The Physical Layer provides hardware capacity. The Virtualization Layer manage hardware abstraction into virtual machines. The Containerization Layer hosts microservice workloads. Finally, the Monitoring Layer observes telemetry to drive orchestration logic.

\section{Physical and Network Foundation}

The foundation rests on failure domain isolation, where hardware is segmented to prevent single points of failure. The network architecture implements a three-segment isolation model:
\begin{itemize}
    \item \textbf{Workload Segments:} Isolated environments for distinct microservice groups.
    \item \textbf{Observability Segment:} A dedicated segment for telemetry and administrative traffic.
    \item \textbf{Management Segment:} Restricted access for orchestration agents.
\end{itemize}

Routing follows a default-deny posture, where inter-segment communication requires explicit authorization policies.

\section{Virtualization and VM Placement}

\subsection{Mini-Cluster Partitioning}
To optimize scheduling, the virtual fleet is logically partitioned into mini-clusters. This strategy confines the search space of placement algorithms, reducing computational complexity while ensuring microservice co-location for high-frequency communication patterns.

\subsection{Hybrid WOA-DE Placement Algorithm}
The mapping of Virtual Machines to Physical Hosts is determined by a hybrid algorithm combining Whale Optimisation (WOA) and Differential Evolution (DE). This combination ensures both robust global exploration and precise local exploitation.

\begin{algorithm}[H]
\DontPrintSemicolon
\KwIn{VM Request $\vec{R}$, Host Metrics $M$, Population $N$}
\KwOut{Optimal Mapping $H^*$}
Initialise population of candidate host assignments\;
\While{Termination criteria not met}{
    Evaluate Fitness $f(X)$ using weighted Euclidean distance\;
    \tcp{WOA Update Phase}
    Update positions via shrinking encirclement and spiral movement\;
    \tcp{DE Update Phase}
    Apply mutation and crossover to maintain population diversity\;
    Select survivors through tournament selection\;
    Update global best $H^*$\;
}
\Return $H^*$\;
\caption{Hybrid WOA-DE Provisioning Algorithm}
\label{alg:woa_de}
\end{algorithm}

\subsection{Fitness Function}
Placement quality is evaluated using a multi-objective weighted Euclidean distance formula. The objective is to minimize the resource differential between the request and the host's available capacity.

\begin{equation}
\text{Fitness} = \sqrt{w_C (\Delta C)^2 + w_R (\Delta R)^2 + w_I (\Delta I)^2 + w_E (\Delta E)^2}
\end{equation}

Where $\Delta C, \Delta R, \Delta I, \Delta E$ represent standardized Z-scores for CPU, RAM, I/O, and Energy respectively, subject to $\sum w_i = 1$.

\section{Intelligence and Security}

\subsection{Dual-Agent Decision System}
Autonomous remediation is managed by two cognitive agents. The \textbf{Planner} ingest telemetry and formulates strategies. The \textbf{Critic} serves as a safety validator, evaluating proposals against physical headroom and safety policies.

\subsection{Risk-Based Escalation Matrix}
Actions are classified by risk to determine the authorization level. Table \ref{tab:risk_matrix} defines the escalation protocols.

\begin{table}[htbp]
\centering
\caption{Risk-Based Escalation Matrix}
\label{tab:risk_matrix}
\begin{tabular}{@{}lll@{}}
\toprule
\textbf{Risk Level} & \textbf{Action Type} & \textbf{Protocol} \\ \midrule
Low & Scaling Up, Cloning & Fully Autonomous \\
Medium & Live Migration, Scaling Down & Autonomous with Audit Log \\
High & Stop/Delete Resource, Drain Node & Manual Admin Approval (Email) \\ \bottomrule
\end{tabular}
\end{table}

\chapter{Implementation}

%% ---------------------------------------------------------------
%% TODO (student): Complete the following before final submission:
%%   - Replace all [PLACEHOLDER] figures with actual screenshots
%%   - Fill in the VM inventory table from Proxmox
%%   - Fill in the Kubernetes object table from kubectl output
%%   - Add performance measurement results if available
%% ---------------------------------------------------------------

\section{Introduction}

This chapter details the practical realization of the architecture designed in Chapter 2. It identifies the chosen technologies and grounds the theoretical layers in a concrete implementation hosting a distributed medical information system.

\section{Technology Selection and Justification}

The implementation leverages professional open-source tools. Table \ref{tab:tech_stack} summarizes the mapping between the architectural layers and the selected technology stack.

\begin{table}[htbp]
\centering
\caption{Implementation Technology Mapping}
\label{tab:tech_stack}
\begin{tabular}{@{}lll@{}}
\toprule
\textbf{Layer} & \textbf{Chosen Technology} & \textbf{Justification} \\ \midrule
Virtualization & Proxmox VE (KVM/LXC) & Robust API, hardware efficiency \\
Containerization & Kubernetes & Industry standard for microservices \\
Monitoring & Prometheus & Pull-based metrics, native alerting \\
Network & OPNsense & State-of-the-art virtual routing \\ \bottomrule
\end{tabular}
\end{table}

\section{Implementation Environment}

The physical cluster comprises 8 nodes (\textit{pfe0} to \textit{pfe7}). Table \ref{tab:physical_hosts} provides the inventory of the infrastructure substrate.

\begin{table}[htbp]
\centering
\caption{Physical Host Inventory}
\label{tab:physical_hosts}
\begin{tabular}{@{}llllll@{}}
\toprule
\textbf{Host} & \textbf{Role} & \textbf{CPU} & \textbf{RAM} & \textbf{Storage} & \textbf{OS} \\ \midrule
pfe0-pfe7 & Hypervisor & [TBC] & [TBC] & [TBC] & Proxmox VE \\ \bottomrule
\end{tabular}
\end{table}

\section{Virtualization with Proxmox VE}

Proxmox VE implements the virtualization layer, managing the lifecycle of VMs through the QEMU/KVM hypervisor. Resource pools are configured to isolate tenant workloads. Table \ref{tab:vm_inv} provides a preliminary VM inventory.

\begin{table}[htbp]
\centering
\caption{Virtual Machine Inventory}
\label{tab:vm_inv}
\begin{tabular}{@{}lllll@{}}
\toprule
\textbf{VM Name} & \textbf{vCPU} & \textbf{RAM} & \textbf{Service Group} & \textbf{Mini-Cluster} \\ \midrule
k8s-master & [TBC] & [TBC] & Management & N/A \\
k8s-worker-1 & [TBC] & [TBC] & Hospital & mini-cluster-1 \\
k8s-worker-2 & [TBC] & [TBC] & Ministry & mini-cluster-2 \\ \bottomrule
\end{tabular}
\end{table}

\begin{figure}[htbp]
\centering
\fbox{\rule{0pt}{5cm}\rule{0.9\linewidth}{0pt}}
\caption{[PLACEHOLDER — Screenshot of the Proxmox VE Management Dashboard]}
\label{fig:proxmox_dashboard}
\end{figure}

\section{Container Orchestration with Kubernetes}

Kubernetes manages the microservices deployment. The Chapter 2 mini-cluster logic is realized through Kubernetes Namespaces.

\begin{table}[htbp]
\centering
\caption{Kubernetes Namespace Mapping}
\label{tab:k8s_ns}
\begin{tabular}{@{}lll@{}}
\toprule
\textbf{Namespace} & \textbf{Service Group} & \textbf{Network Segment} \\ \midrule
hospital & Hospital MIS & VLAN 10 (Hospital) \\
ministry & Ministry MIS & VLAN 20 (Ministry) \\
monitoring & Observability & VLAN 30 (Monitor) \\ \bottomrule
\end{tabular}
\end{table}

\section{Deployed Application — Medical Information System}

The system hosts two major microservice groups. The communication patterns between these groups are managed through explicit firewall rules in the OPNsense gateway.

\subsection{Ministry Microservices}
These services handle national registries and security. Based on the implementation, the group includes: \textbf{auth} (Authentication), \textbf{blockchain-service} (Data integrity), \textbf{cpabe-service} (Encryption), \textbf{rabbitmq} (Messaging), and \textbf{frontend-ministry}.

\subsection{Hospital Microservices}
These services manage local clinical operations. The group includes: \textbf{patients} (Records), \textbf{rendez-vous} (Appointments), \textbf{pharmacy} (Inventory), \textbf{billing} (Interventions), and \textbf{frontend-hospital}.

\begin{figure}[htbp]
\centering
\fbox{\rule{0pt}{5cm}\rule{0.9\linewidth}{0pt}}
\caption{[PLACEHOLDER — Diagram of the Application Microservices Architecture]}
\label{fig:app_arch}
\end{figure}

\section{Monitoring and Network Implementation}

\subsection{Prometheus Telemetry}
The Prometheus server is configured to scrape metrics from every node. Key targets include \texttt{node\_exporter} (Port 9100) and \texttt{scaphandre} (Port 8080). Alert rules are defined for resource breaches.

\subsection{Overlay Network Segmentation}
The OPNsense virtual router implements three logical interfaces. Routing is restricted to prevent unauthorized lateral movement between the Hospital and Ministry segments.
\begin{itemize}
    \item \textbf{Hospital Segment:} Subnet 10.10.10.0/24.
    \item \textbf{Ministry Segment:} Subnet 10.20.20.0/24.
    \item \textbf{Monitoring Segment:} Subnet 10.30.30.0/24.
\end{itemize}

\begin{figure}[htbp]
\centering
\fbox{\rule{0pt}{5cm}\rule{0.9\linewidth}{0pt}}
\caption{[PLACEHOLDER — Screenshot of the OPNsense Firewall Rules]}
\label{fig:firewall_config}
\end{figure}

\section{Results and Validation}

[PLACEHOLDER — Detailed analysis of algorithm efficiency and scaling performance logs].

\begin{figure}[htbp]
\centering
\fbox{\rule{0pt}{5cm}\rule{0.9\linewidth}{0pt}}
\caption{[PLACEHOLDER — Screenshot of a scaling event triggered by the Planner agent]}
\label{fig:scaling_log}
\end{figure}

\chapter*{General Conclusion}
\addcontentsline{toc}{chapter}{General Conclusion}

This thesis has demonstrated the feasibility of an autonomous provisioning solution for high-load web applications. By utilizing a hierarchical architecture and hybrid meta-heuristic algorithms, the system effectively optimizes resource utilization while maintaining service availability. The implementation of a medical information system on a distributed Proxmox/Kubernetes cluster validated the theoretical framework. The integration of risk-stratified agents ensures that automated remediation remains safe and forensic auditability is preserved.

\begin{thebibliography}{99}
\bibitem{balabanov2025dcim}
G.~Balabanov and N.~S. Apostolov, ``Data Centre Infrastructure Monitoring,'' in \textit{60th International Scientific Conference on ICEST}, IEEE, 2025.
\bibitem{mirjalili2016woa}
S.~Mirjalili and A.~Lewis, ``The Whale Optimisation Algorithm,'' \textit{Advances in Engineering Software}, vol. 95, pp. 51--67, 2016.
\bibitem{storn1997de}
R.~Storn and K.~Price, ``Differential Evolution --- A Simple and Efficient Heuristic for Global Optimization over Continuous Spaces,'' \textit{Journal of Global Optimization}, vol. 11, no. 4, pp. 341--359, 1997.
\bibitem{lai2025energy}
X.~Lai, X.~Zheng, and W.~Lu, ``Energy-aware Virtual Machine Migration based on Dual Thresholds,'' in \textit{2025 IEEE 2nd International Conference on ICBDSE}, IEEE, 2025.
\bibitem{zhang2020intelligent}
P.~Zhang, M.~Zhou, and X.~Wang, ``An Intelligent Optimization Method for Optimal Virtual Machine Allocation,'' \textit{IEEE Transactions on Automation Science and Engineering}, vol. 17, no. 4, pp. 1725--1735, 2020.
\bibitem{saxena2022secure}
D.~Saxena, I.~Gupta, et al., ``A Secure and Multiobjective Virtual Machine Placement Framework,'' \textit{IEEE Systems Journal}, vol. 16, no. 2, pp. 3163--3174, 2022.
\end{thebibliography}

\end{document}
Modern software development has been shaped by evolving scalability demands and the need for efficient resource utilization. Virtualization and containerization have fundamentally altered how applications are developed, deployed, and managed. This chapter establishes the theoretical foundation for these paradigm shifts, covering their core concepts, operational characteristics, and relationship to microservices architecture, before critically examining whether organizations face a genuine imperative to adopt them.
>>>>>>> Stashed changes

% ================================================================
\section{Virtualization and Containerization}
% ================================================================

\subsection{Virtualization}

Virtualization enables the creation of simulated computing environments abstracted
from physical hardware. A hypervisor partitions physical resources into isolated
virtual machines (VMs), each with its own allocated CPU, memory, storage, and
network interfaces \cite{popek1974formal}. By decoupling software execution from
hardware, multiple operating systems can coexist on a single server while remaining
operationally independent \cite{rosenblum2005virtual}.

Key advantages include resource consolidation, reduced infrastructure costs, live
migration of running VMs between hosts, rapid provisioning of environments, and
flexible workload placement \cite{vmware2006consolidation, barham2003xen}.

% Figure 1.1 – placed directly after the paragraph describing virtualisation.
\begin{figure}[H]
    \centering
    \includegraphics[width=0.75\textwidth]{graphics/vm_architecture.png}
    \caption{Virtual Machine Architecture Overview}
    \label{fig:vm_architecture}
\end{figure}
\FloatBarrier

\subsection{Containerization}

Containerization is a lightweight alternative that operates at the operating-system
level rather than the hardware level. Containers encapsulate application code with
its dependencies into portable units that share the host OS kernel, using kernel
features such as \textit{namespaces} for isolation and \textit{control groups} for
resource limiting \cite{merkel2014docker, soltesz2007container}.

Key advantages include faster startup times, higher deployment density, consistent
behavior across environments, natural alignment with microservices and CI/CD
pipelines, and declarative version-controlled definitions that simplify rollbacks
\cite{bernstein2014containers, pahl2015containerization}.

\FloatBarrier

% ================================================================
\section{Virtualized and Containerized Environments}
% ================================================================

\subsection{Virtual Machine Architecture}

Hypervisors come in two types: Type~1 (bare-metal) run directly on hardware, while
Type~2 run atop a host OS \cite{rosenblum2005virtual}. Each VM includes a full
guest OS and virtualized hardware managed by the hypervisor, providing strong
isolation at the cost of higher resource overhead and longer boot times
\cite{barham2003xen}.

\subsection{Container-Based Environments}

Container runtimes such as Docker or containerd interface with the host kernel to
create isolated execution contexts. Namespaces prevent containers from observing
processes outside their scope; control groups prevent resource monopolization; and
union file systems allow containers to share common base layers while maintaining
isolated writable layers \cite{merkel2014docker, soltesz2007container}.

Container orchestration platforms such as Kubernetes add automated scheduling,
service discovery, health monitoring, and rolling updates to manage containers at
scale \cite{burns2016borg}.

% Figure 1.2 – placed directly after the container-based environments description,
% pinned here by [H] + \FloatBarrier so it cannot drift into later sections.
\begin{figure}[H]
    \centering
    \includegraphics[width=0.75\textwidth]{graphics/container_environment.png}
    \caption{Container-Based Environment and Components}
    \label{fig:container_environment}
\end{figure}
\FloatBarrier

% ================================================================
\section{Microservices Architecture}
% ================================================================

\subsection{Definition and Principles}

Microservices structures applications as collections of loosely coupled,
independently deployable services, each implementing a specific business capability
and owning its own data store \cite{newman2015microservices}. Core principles
include single responsibility, bounded contexts from domain-driven design,
lightweight communication via HTTP APIs or message queues, decentralized governance,
and built-in resilience patterns such as circuit breakers and timeouts
\cite{fowler2014microservices, richardson2018microservices}.

\subsection{Benefits and Challenges}

Microservices enable independent deployability, technology heterogeneity, granular
scalability, team autonomy, and incremental modernization of legacy systems
\cite{newman2015microservices, fowler2014microservices}.

However, they introduce significant complexity: distributed debugging requires
sophisticated monitoring and tracing; cross-service transactions demand eventual
consistency models; network latency introduces new failure points; and substantial
operational infrastructure --- service discovery, load balancing, orchestration ---
becomes necessary \cite{richardson2018microservices}.

% ================================================================
\section{Microservices versus Monolithic Architecture}
% ================================================================

Monolithic architecture integrates all components --- UI, business logic, and data
access --- into a single codebase sharing a common database, deployed as one unit
\cite{fowler2014microservices}. Components communicate through in-process calls,
simplifying initial development and ensuring strong data consistency.

The key structural differences are summarized as follows. Monoliths exhibit tight
coupling where changes ripple across the codebase, while microservices enforce loose
coupling through stable interfaces. Monoliths use a shared database with strong
consistency; microservices use per-service databases, complicating cross-service
transactions. Monoliths scale by replicating the entire application; microservices
enable selective scaling of individual components. Failures in a monolith can
cascade system-wide, whereas microservices contain failures within service
boundaries --- though distributed failures introduce new failure modes
\cite{newman2015microservices, richardson2018microservices, fowler2014microservices}.

% ================================================================
\section{Critical Analysis: The Imperative for Microservices with Containerization
and Virtualization}
% ================================================================

\subsection{Limitations of Traditional Deployment}

Traditional deployment on physical or dedicated virtual servers suffers from
documented inefficiencies. Infrastructure often operates below 20--30 percent
average utilization while being over-provisioned for peak loads, representing
substantial capital waste \cite{vmware2006consolidation}. Manual provisioning is
slow and inconsistent, limiting organizational responsiveness. Environmental
discrepancies between development and production cause deployment failures
\cite{merkel2014docker}, and monolithic architectures offer limited fault tolerance
\cite{rosenblum2005virtual}.

\subsection{Should Organizations Transition to Microservices and
Virtualization/Containerization Environments?}

The strategic question is straightforward to pose but difficult to answer:
\textit{should organizations transition to microservices supported by virtualization
and containerization?}

The case for transition is clear in principle. Containerization delivers consistency
and resource efficiency; virtualization enables consolidation; microservices unlock
deployment velocity and team autonomy \cite{merkel2014docker, richardson2018microservices}.
Yet the associated costs are substantial. Container orchestration demands specialized
expertise, migration requires significant investment in infrastructure and personnel,
and the operational maturity required for distributed microservices exceeds what most
organizations have built through traditional models \cite{burns2016borg, bernstein2014containers}.

The decision is further complicated by inherent contradictions: larger organizations
may benefit most from consolidation yet face more complex migrations; stable legacy
systems may gain little from disruption; and containerization without architectural
transformation yields limited benefits, while microservices without containerization
introduce different operational burdens. Organizations must determine not only
\textit{whether} to adopt these technologies but \textit{in what sequence and scope}
--- questions for which industry guidance remains conflicting and empirical evidence
limited \cite{barham2003xen}.

\FloatBarrier

% ============================================================
% CHAPTER 2 – Conception and Architectural Design
% ============================================================
% ============================================================
% CHAPTER 2 – Conception and Architectural Design
% ============================================================
\chapter{Conception and Architectural Design}
\label{chap:conception}

\section{Introduction}
\label{sec:ch2_intro}

This chapter presents the full architectural design of a cluster infrastructure for
a high-demand web application, organised around three pillars: a proof of concept
defining clear validation criteria before any large-scale deployment, a four-layer
infrastructure stack from physical hardware to a monitoring plane, and an intelligence
and security layer combining a natural-language assistant, a predictive model, and a
zero-trust security design. The three pillars are mutually dependent: the monitoring
layer feeds the placement algorithm, the security layer protects that data, and the
predictive model learns from monitoring history.

% ================================================================
\section{Proof of Concept: Goals and Validation}
\label{sec:poc}
% ================================================================

\subsection{Design Objectives}
\label{subsec:poc_objectives}

The system faces two requirements that pull in opposite directions:
\textit{resource efficiency} --- packing workloads tightly to maximise hardware
utilisation --- and \textit{continuous availability} --- keeping every service
reachable despite hardware failures or traffic spikes. The architecture resolves
this tension through intelligent scheduling and scaling that adapts in real time.
The application is built as loosely coupled microservices so that each service scales
independently and a fault in one causes graceful degradation rather than a full
outage \cite{newman2015microservices}.

\subsection{Validation Criteria}
\label{subsec:poc_criteria}

A time-bounded proof-of-concept phase verifies that all core mechanisms work together
before full-scale deployment. All four conditions must hold simultaneously:

\begin{itemize}
    \item \textbf{Automated provisioning:} a new VM joins the cluster automatically,
    within a defined time bound, with no manual steps.
    \item \textbf{Responsive scaling:} injected anomalies (CPU saturation, memory
    pressure, high error rates, packet loss) trigger the correct remediation within
    one monitoring cycle.
    \item \textbf{Constraint-respecting scheduling:} containers are placed without
    exceeding VM resource limits and rebalance automatically when VMs change.
    \item \textbf{Preserved availability:} the application remains reachable
    throughout all provisioning, scaling, and rebalancing operations.
\end{itemize}

% ================================================================
\section{Infrastructure Layers}
\label{sec:layers}
% ================================================================

The cluster is organised into four superimposed layers. The physical layer provides
raw compute, storage, and network capacity. The virtualisation layer abstracts it
into managed virtual machines. The containerisation layer runs microservices as
containers across those VMs. The monitoring layer observes all three and drives
automated responses. Figure~\ref{fig:layers} illustrates this stack and the data
flows between layers.

% Figure 2.1 – pinned immediately after the section intro paragraph that references it.
% \FloatBarrier stops it from drifting into subsection 2.3.1.
\begin{figure}[H]
    \centering
    \vspace{4pt}
    \includegraphics[width=0.80\textwidth]{layers.pdf}
    \caption{The four-layer infrastructure stack. Each layer builds on the one
    below it; the monitoring and intelligence layer (L4) observes all others and
    drives automated responses.}
    \label{fig:layers}
    \vspace{4pt}
\end{figure}
\FloatBarrier

% ----------------------------------------------------------------
\subsection{Physical and Network Foundation}
\label{subsec:physical}
% ----------------------------------------------------------------

Physical servers are connected through a managed switching fabric; hardware firewalls
at the network edge enforce ingress and egress policies before traffic reaches
internal services. The layer is built around explicit \textit{failure domains}:
replicas are spread across at least two distinct domains so that no single hardware
event eliminates all instances of a service. All network equipment is deployed in
redundant pairs \cite{barroso2013datacenter}.

\FloatBarrier

% ----------------------------------------------------------------
\subsection{Virtualisation Layer: Placement, Partitioning, and Scaling}
\label{subsec:virtual}
% ----------------------------------------------------------------

\paragraph{Initial Sizing and Mini-Cluster Organisation.}
Each new VM receives an initial resource allocation derived from a workload profile
--- benchmarks, service-level objectives, or historical data from similar services.
A set of standard \textbf{size classes} is defined on top of this baseline, providing
pre-validated resource ratios for the most common workload types and simplifying
provisioning decisions for well-characterised services.
The monitoring layer tracks usage continuously and triggers vertical scaling if
sustained demand exceeds the initial allocation.

The VM fleet is partitioned into $k$ logical \textit{mini-clusters}, each grouping
VMs whose workloads communicate frequently. This reduces scheduling scope from $n$
physical machines to $n/k$ per decision (a factor-$k$ complexity reduction), keeps
related traffic local, and prevents hotspots \cite{zhang2020intelligent}.
Figure~\ref{fig:mini_clusters} illustrates the partitioning. The value of $k$ is
either set statically by the application designer, or inferred dynamically by a
machine learning model observing inter-service communication at runtime.

% Figure 2.2 – pinned after the mini-cluster paragraph that references it.
% \FloatBarrier keeps it above the Provisioning Algorithm paragraph.
\begin{figure}[H]
    \centering
    \vspace{4pt}
    \includegraphics[width=0.68\textwidth]{clusters.pdf}
    \caption{VM fleet mini-cluster partitioning. Scheduling scope is reduced by a
    factor of $k$; intra-cluster traffic is maximised and cross-cluster traffic
    minimised.}
    \label{fig:mini_clusters}
    \vspace{4pt}
\end{figure}
\FloatBarrier

\paragraph{Provisioning Algorithm and Fitness Evaluation.}
VM placement uses a hybrid of the Whale Optimisation Algorithm (WOA)
\cite{mirjalili2016woa} and Differential Evolution (DE) \cite{storn1997de}.
At each iteration the candidate population is split into two disjoint subsets.
The \textbf{WOA subset} applies encirclement, spiral, or random-agent position
updates; the \textbf{DE subset} applies mutation, crossover, and greedy selection.
The global best $X^*$ is refreshed after both subsets. This prevents premature
convergence (a weakness of pure WOA) and the aimless drift that arises in DE when
population diversity collapses.

The fitness function scores a candidate host against the VM's request across four
dimensions (Table~\ref{tab:fitness_dims}). Each metric is first z-score normalised:

\begin{equation}
    X_i^{*} = \frac{X_i - \bar{X}}{\sigma_X}
    \label{eq:zscore}
\end{equation}

and combined into a weighted Euclidean distance:

\begin{equation}
    \text{fitness} = \sqrt{\,w_1(\Delta C)^2 + w_2(\Delta R)^2
        + w_3(\Delta IO)^2 + w_4(\Delta E)^2},\quad \textstyle\sum w_i = 1
    \label{eq:fitness}
\end{equation}

where $\Delta X$ is the gap between what the VM requests and what the host offers.
A lower score indicates a better placement.

\begin{table}[H]
    \centering
    \caption{Fitness Function Resource Dimensions}
    \label{tab:fitness_dims}
    \renewcommand{\arraystretch}{1.4}
    \begin{tabular}{|l|l|l|l|}
        \hline
        \textbf{Dimension} & \textbf{Symbol} & \textbf{Unit} & \textbf{Rationale} \\
        \hline\hline
        CPU Capacity   & $C$  & vCores & Prevents compute starvation \\
        \hline
        RAM            & $R$  & GB     & Avoids memory-induced swapping \\
        \hline
        I/O Throughput & $IO$ & MB/s   & Guards storage and network bandwidth \\
        \hline
        Energy Draw    & $E$  & J/hr   & Minimises operational energy cost \\
        \hline
    \end{tabular}
\end{table}

The full iterative procedure is shown in Figure~\ref{fig:woa_de} and summarised
below.

\textit{Initialisation.} Sample $n$ random VM-to-host mappings; set $X^*$ to the
lowest-fitness candidate; set $a_0{=}2$, scaling factor $F$, crossover rate
$CR{=}0.5$.

\textit{Main loop} (for $t = 0,\ldots,T-1$):

\begin{enumerate}
    \item \textbf{Parameter update.}
    \begin{equation}
        a = 2\!\left(1-\tfrac{t}{T}\right),\quad
        A = 2ar_1 - a,\quad C = 2r_2,\quad p\sim\mathrm{U}(0,1)
        \label{eq:woa_params}
    \end{equation}

    \item \textbf{WOA update} (for each $X_i$ in WOA subset):
    \begin{equation}
        X_i(t{+}1) =
        \begin{cases}
            X^* - A\cdot|CX^* - X_i|
                & p < 0.5,\;|A| < 1\;\text{(exploitation)}\\[4pt]
            X_{\mathrm{rand}} - A\cdot|CX_{\mathrm{rand}} - X_i|
                & p < 0.5,\;|A| \ge 1\;\text{(exploration)}\\[4pt]
            |X^*-X_i|\cdot e^{bl}\cos(2\pi l)+X^*
                & p \ge 0.5\;\text{(spiral)}
        \end{cases}
        \label{eq:woa_update}
    \end{equation}
    Refresh $X^*$ from the WOA subset.

    \item \textbf{DE update} (for each $X_j$ in DE subset):
    \begin{align}
        V_j   &= X^* + F\,(X_{r_1} - X_{r_2})
            \label{eq:de_mut}\\[2pt]
        U_j[k]&= \begin{cases}V_j[k]&\mathrm{rand}()\le CR\\
                              X_j[k]&\text{otherwise}\end{cases}
            \label{eq:de_cx}\\[2pt]
        X_j   &\leftarrow U_j\;\text{ if }f(U_j)<f(X_j)
            \label{eq:de_sel}
    \end{align}
    Refresh $X^*$ from both subsets.
\end{enumerate}

\textit{Output}: $X^*$ --- the recommended VM-to-host placement for the mini-cluster.

% Figure 2.3 – pinned after the algorithm description, before the scaling paragraph.
\begin{figure}[H]
    \centering
    \vspace{4pt}
    \includegraphics[width=\textwidth]{VM-Provisioning.pdf}
    \caption{Hybrid WOA--DE provisioning algorithm (swimlane view). Left: WOA
    update phase. Right: DE mutation, crossover, and selection. Centre: shared
    $X^*$ synchronisation.}
    \label{fig:woa_de}
    \vspace{4pt}
\end{figure}
\FloatBarrier

% --- NEW: Algorithm properties and limitations ---
The hybrid design exhibits five desirable properties.
(i)~\textbf{Dual exploration:} WOA spiral and random-agent modes explore the
solution space from directions anchored to $X^*$ and to random population members,
while DE mutation explores directions defined by the algebraic difference
$X_{r_1}-X_{r_2}$; the two mechanisms are geometrically independent.
(ii)~\textbf{Preservation of the global best:} both WOA contraction and DE
base-vector choice orient the population toward $X^*$, which is refreshed after
each phase so that improvements found during the WOA phase immediately influence
the DE phase.
(iii)~\textbf{Greedy selection:} DE discards any trial mapping that does not
improve upon its parent, ensuring monotonically non-increasing population fitness
across iterations.
(iv)~\textbf{Parameter sensitivity reduction:} mixing both paradigms distributes
the optimisation workload across two parameter sets whose failure modes are
uncorrelated.
(v)~\textbf{Suitability for discrete placement spaces:} the fitness function
handles the required continuous-to-discrete decoding, so both update rules apply
without modification.
The principal limitations are the O($n$) fitness evaluations per iteration, which
becomes significant at large fleet sizes, and the discrete decoding step required
before each evaluation.
% --- END NEW ---

\paragraph{Runtime Scaling Mechanisms.}
\textbf{Horizontal scaling} creates a new VM in the same mini-cluster when request
volume exceeds what a single instance can handle. \textbf{Vertical scaling} adds
vCPUs, memory, or storage when a specific resource is saturated; if the host lacks
headroom, the VM is live-migrated within the mini-cluster \cite{lai2025energy}.
All trigger conditions and actions are listed in Table~\ref{tab:scaling}.

\begin{table}[H]
    \centering
    \caption{Scaling Trigger Conditions and Remediation Actions}
    \label{tab:scaling}
    \renewcommand{\arraystretch}{1.4}
    \begin{tabular}{|p{2.1cm}|p{4.0cm}|p{1.9cm}|p{4.8cm}|}
        \hline
        \textbf{Type} & \textbf{Trigger} & \textbf{Severity} & \textbf{Action} \\
        \hline\hline
        Horizontal & HTTP 5xx $> 1\%$        & Warning     & Spawn new VM in mini-cluster \\
        \hline
        Horizontal & HTTP 5xx $> 5\%$        & Critical    & Spawn new VM immediately \\
        \hline
        Horizontal & Packet loss $> 2\%$      & Critical    & Spawn new VM, rebalance pool \\
        \hline
        Vertical   & CPU $> 80\%$            & Warning     & $+1$ vCPU \\
        \hline
        Vertical   & CPU $> 95\%$            & Critical    & $+1$ vCPU; migrate if no headroom \\
        \hline
        Vertical   & RAM $> 60\%$            & Warning     & $+20\%$ memory \\
        \hline
        Vertical   & RAM $> 80\%$            & Critical    & $+20\%$ memory; migrate if needed \\
        \hline
        Vertical   & Disk I/O free $< 20\%$  & Warn./Crit. & $+15$\,GB storage \\
        \hline
    \end{tabular}
\end{table}

\paragraph{Virtual Network Architecture.}
An overlay network decouples logical topology from physical cabling, enforcing
segment boundaries in software on each hypervisor node without reconfiguring
physical equipment. This approach provides greater flexibility and portability than
VLAN-based segmentation, and is particularly suited to environments where the
underlying infrastructure cannot be modified.
Three logical segments are defined: primary tenant workloads, a second isolated
environment (e.g.\ staging), and a dedicated monitoring segment. All inter-segment
traffic routes through a central virtual router enforcing \textit{default-deny}:
unmatched packets are silently dropped and cross-segment flows require explicit
bidirectional rules. The monitoring segment has targeted read-only access to other
segments; no reverse flow is permitted. An internal name-resolution service
eliminates reliance on external public resolvers, removing the DNS-redirection risk.

\FloatBarrier

% ----------------------------------------------------------------
\subsection{Container Scheduling and Placement}
\label{subsec:containers}
% ----------------------------------------------------------------

Each microservice runs in its own container, isolated by kernel namespaces and
control groups. Containers are assigned to VMs using a \textbf{Best Fit} strategy:
the scheduler selects the VM with the least remaining capacity that can still meet
the container's requirements, filling VMs before touching less-loaded ones. This
minimises the number of active VMs and therefore total energy and licensing costs
\cite{coffman1984bin}. Mini-cluster boundaries are respected: containers from the
same mini-cluster are preferentially co-located on VMs in that mini-cluster to
preserve network locality.

\FloatBarrier

% ----------------------------------------------------------------
\subsection{Observability and Automated Response}
\label{subsec:monitoring}
% ----------------------------------------------------------------

The monitoring layer consists of a \textbf{telemetry subsystem} and an \textbf{alert
manager}. The telemetry subsystem collects metrics from all layers using a pull-based
model \cite{nguyen2022loadbalancing}, spanning all observable infrastructure
categories: physical servers (CPU utilisation, memory consumption, disk I/O
throughput, and energy draw), virtual machines (per-VM CPU steal time and
working-set size), and physical and virtual network devices (port utilisation,
packet loss counters, inter-VM traffic volumes).
% --- NEW: retention and correlation ---
Data is stored with \textbf{time-windowed retention policies} that preserve
high-resolution data for anomalous intervals and apply progressive downsampling to
normal ones. Co-located metrics --- such as the CPU and memory utilisation of all
VMs sharing the same physical host --- are stored with their correlational
relationships preserved, enabling \textbf{correlation-based anomaly detection}
rather than isolated threshold crossings.
% --- END NEW ---

The alert manager maps structured events to actions: scaling requests to the
virtualisation layer, or operator notifications for situations requiring human review.
Grouping and inhibition rules prevent alert storms during cascading failures by
consolidating related events and suppressing symptoms when a root-cause alert is
already active.

\FloatBarrier

% ================================================================
\section{Intelligence and Security}
\label{sec:intelligence}
% ================================================================

Static thresholds and manual oversight are insufficient at scale. The architecture
adds three intelligent components --- a conversational assistant, a predictive model,
and a two-agent decision system --- all protected by a zero-trust security design.

% ----------------------------------------------------------------
\subsection{Conversational Administration Assistant}
\label{subsec:assistant}
% ----------------------------------------------------------------

A natural-language assistant is embedded in the cluster management interface,
allowing operators to query cluster state in plain language without navigating
dashboards or memorising CLI syntax. Its autonomy is bounded: read queries execute
immediately, but any write action is suspended until the operator explicitly confirms.
The assistant is also the single human-facing interface through which the predictive
model surfaces recommendations.

% ----------------------------------------------------------------
\subsection{Predictive Resource Management}
\label{subsec:predictive}
% ----------------------------------------------------------------

The predictive model operates across two horizons. \textbf{Before VM creation (H1)},
it recommends an initial resource allocation based on service type, expected load,
and usage patterns of similar running services, reducing early scaling events.
\textbf{At runtime (H2)}, it applies time-series forecasting to detect trajectories
heading toward a scaling threshold and sends pre-emptive recommendations through
the assistant for operator approval.
% --- NEW: dynamic thresholds, adaptive weights, incremental learning ---
The model serves three additional functions.
First, \textbf{dynamic alert thresholds}: rather than relying on the fixed values
in Table~\ref{tab:scaling}, it continuously recalibrates warning and critical
thresholds using recent telemetry as a learned baseline, raising them during
high-traffic events and tightening them during quiet periods to reduce false
positives.
Second, \textbf{adaptive fitness weights}: the model continuously identifies the
current structural bottleneck of the cluster (e.g.\ memory pressure or I/O
saturation) and dynamically adjusts the weights $w_1$--$w_4$ in
Equation~\ref{eq:fitness}, so that subsequent placement decisions automatically
prioritise the constrained resource.
Third, to remain effective as workloads evolve, the model employs
\textbf{incremental self-supervised learning}: at regular intervals it ingests the
latest cluster metrics and updates its internal representations using its own recent
predictions as soft labels, adapting to new baseline behaviours without requiring
external ground-truth datasets.
% --- END NEW ---

% ----------------------------------------------------------------
\subsection{Autonomous Decision-Making and Risk-Based Escalation}
\label{subsec:escalation}
% ----------------------------------------------------------------

Remediation is handled by a \textbf{Planner--Critic} system. The \textbf{Planner}
reads telemetry, bottleneck signals, and action history; when it detects a problem
it proposes a concrete action and assigns an initial risk level. The \textbf{Critic}
independently validates the proposal against three conditions --- host headroom,
absence of runaway loops, and alignment with the predicted bottleneck --- and may
approve, reject, or escalate the risk. A critical constraint: the Critic may
\textit{raise} but is structurally prohibited from \textit{lowering} the Planner's
risk classification, guarding against adversarial misclassification.
% --- NEW: orchestration layer isolation ---
The agents do not hold direct credentials to the infrastructure management
interfaces; they submit structured proposals to an independent \textbf{orchestration
layer} that validates the execution context before dispatching commands --- a
compromised agent can propose but cannot execute unilaterally.
% --- END NEW ---
Table~\ref{tab:escalation} shows the routing; Figure~\ref{fig:planner_critic}
illustrates the full pipeline.

\begin{table}[H]
    \centering
    \caption{Risk-Based Escalation Matrix}
    \label{tab:escalation}
    \renewcommand{\arraystretch}{1.5}
    \begin{tabular}{|p{1.7cm}|p{3.1cm}|p{3.7cm}|p{4.1cm}|}
        \hline
        \textbf{Risk} & \textbf{Example Actions} &
        \textbf{System Response} & \textbf{Human Role} \\
        \hline\hline
        Low &
            Add resources, clone instance &
            Auto-executed upon Critic approval &
            Informed via audit log only \\
        \hline
        Medium &
            Vertical scaling, threshold adjustment &
            Auto-executed; reasoning logged &
            Notified via assistant; may override in grace window \\
        \hline
        High &
            VM deletion, forced scale-down &
            Suspended; placed in approval queue with context and rollback plan &
            Must explicitly approve via assistant before execution \\
        \hline
    \end{tabular}
\end{table}

% --- NEW: TTL on approval queue + immutable audit log 3 functions ---
High-risk actions placed in the approval queue are bounded by a
\textbf{time-to-live (TTL)}: an action that does not receive administrator
authorisation within a defined window is automatically rejected and recorded as
expired, preventing an adversary from poisoning the queue and waiting for accidental
approval.
Every event --- Planner proposals, Critic verdicts, execution outcomes, approvals,
rejections, and TTL expirations --- is appended to an \textbf{immutable audit log}
\cite{saxena2022secure}. This log serves three security functions: (i)~forensic
reconstruction of the sequence of events leading to any infrastructure state;
(ii)~detection of gradual privilege-escalation patterns that would be invisible in
any single event; and (iii)~permanent accountability for every infrastructure
change, whether autonomous or human-initiated.
% --- END NEW ---

% Figure 2.4 – pinned after Table 2.3 and before Security Architecture (2.4.4).
% Without float+placeins this figure was drifting to page 17, after the conclusion.
\begin{figure}[H]
    \centering
    \vspace{4pt}
    \includegraphics[width=0.88\textwidth]{pipeline.pdf}
    \caption{Planner--Critic pipeline with risk-based escalation. The Critic may
    raise but never lower the Planner's risk level. High-risk actions require
    explicit operator approval; all outcomes are recorded in an immutable audit log.}
    \label{fig:planner_critic}
    \vspace{4pt}
\end{figure}
\FloatBarrier

% ----------------------------------------------------------------
\subsection{Security Architecture}
\label{subsec:security}
% ----------------------------------------------------------------

Security follows \textit{defence in depth} within a \textit{zero-trust} model:
no entity is trusted by position, every communication is authenticated, and every
component holds the minimum permissions it needs \cite{saxena2022secure}.

\textbf{Network and segmentation.} Hardware firewalls enforce a perimeter allowlist
with bidirectional filtering --- both ingress and egress --- preventing compromised
internal components from establishing unauthorised outbound channels. The
three-segment overlay enforces internal microsegmentation so a compromised workload
cannot move laterally. An anti-DDoS layer absorbs volumetric attacks before they
reach internal services.

\textbf{Identity and integrity.} Every entity holds a cryptographically verifiable
identity; all messages carry a digital signature verified before processing.
Credentials are short-lived and rotated regularly; plaintext channels are prohibited
even on internal segments.

% --- NEW: runtime anomaly detection as dedicated security component ---
\textbf{Runtime anomaly detection.} A dedicated security monitoring agent operates
continuously, subscribing to the same telemetry streams as the alert manager but
applying a behavioural lens oriented toward deviation rather than resource
saturation. It evaluates every observable component against static integrity
baselines and dynamically learned profiles. A \textit{critical} classification
triggers an immediate alert to the administrator, bypassing all aggregation; a
\textit{warning} classification is dispatched at a configured periodic interval,
preserving administrator attention without real-time noise for every minor deviation.
% --- END NEW ---

\textbf{Patch and supply-chain hygiene.} An automated daily process patches every
stack layer from hypervisor to auxiliary services. Components are sourced from
official channels only and verified by signature or content-addressed digest before
installation.

\textbf{Container and VM hardening.} VMs expose only the virtual hardware required;
containers receive only the kernel capabilities they need and filesystems are mounted
read-only wherever possible. Monitoring components hold no provisioning credentials.

\textbf{Cognitive agent security.} The assistant requires administrator authentication
before any query reaches the Planner. Agents submit structured proposals to an
independent orchestration layer that validates execution context before dispatching
commands --- a compromised agent can propose but cannot execute unilaterally. The
Critic's monotonic risk constraint provides a second structural defence.

% ================================================================
\section{Conclusion}
\label{sec:ch2_conclusion}
% ================================================================

The architecture resolves the efficiency--availability tension through intelligent
automation at every layer. The hybrid WOA--DE algorithm minimises resource mismatch
in VM placement; Best Fit scheduling maximises VM utilisation; reactive and predictive
scaling --- augmented by adaptive thresholds and dynamic fitness weights --- keep
capacity matched to demand; the three-segment overlay enforces tenant isolation; and
the Planner--Critic pipeline ensures autonomous actions are proportionate and
auditable. The immutable audit log, digital signatures, TTL-bounded approval queue,
and dedicated security monitoring agent form a coherent defence-in-depth posture
protecting the entire stack. The conversational assistant ties all three intelligent
components into a single human-readable interface that keeps operators in control
without slowing routine operations. Chapter~3 covers the implementation: tool choices,
configuration, and deployment steps.

% ============================================================
% UNIFIED BIBLIOGRAPHY
% ============================================================
\begin{thebibliography}{99}

% --- Chapter 1 ---

\bibitem{popek1974formal}
G.~J.~Popek and R.~P.~Goldberg,
``Formal Requirements for Virtualizable Third Generation Architectures,''
\textit{Communications of the ACM}, vol.~17, no.~7, pp.~412--421, 1974.

\bibitem{rosenblum2005virtual}
M.~Rosenblum and T.~Garfinkel,
``Virtual Machine Monitors: Current Technology and Future Trends,''
\textit{Computer}, vol.~38, no.~5, pp.~39--47, 2005.

\bibitem{vmware2006consolidation}
VMware,
\textit{Server Consolidation and Containment with VMware Infrastructure}.
VMware White Paper, 2006.

\bibitem{barham2003xen}
P.~Barham et al.,
``Xen and the Art of Virtualization,''
in \textit{Proc. ACM SOSP}, Bolton Landing, NY, 2003, pp.~164--177.

\bibitem{merkel2014docker}
D.~Merkel,
``Docker: Lightweight Linux Containers for Consistent Development and Deployment,''
\textit{Linux Journal}, no.~239, 2014.

\bibitem{soltesz2007container}
S.~Soltesz et al.,
``Container-Based Operating System Virtualization: A Scalable, High-Performance
Alternative to Hypervisors,''
in \textit{Proc. ACM EuroSys}, Lisbon, 2007, pp.~275--287.

\bibitem{bernstein2014containers}
D.~Bernstein,
``Containers and Cloud: From LXC to Docker to Kubernetes,''
\textit{IEEE Cloud Computing}, vol.~1, no.~3, pp.~81--84, 2014.

\bibitem{pahl2015containerization}
C.~Pahl,
``Containerization and the PaaS Cloud,''
\textit{IEEE Cloud Computing}, vol.~2, no.~3, pp.~24--31, 2015.

\bibitem{burns2016borg}
B.~Burns, B.~Grant, D.~Oppenheimer, E.~Brewer, and J.~Wilkes,
``Borg, Omega, and Kubernetes,''
\textit{ACM Queue}, vol.~14, pp.~70--93, 2016.

\bibitem{fowler2014microservices}
M.~Fowler and J.~Lewis,
``Microservices,'' martinfowler.com, Mar.\ 2014.
[Online]. Available: \url{https://martinfowler.com/articles/microservices.html}

\bibitem{richardson2018microservices}
C.~Richardson,
\textit{Microservices Patterns}. Manning Publications, 2018.

% --- Chapter 2 ---

\bibitem{newman2015microservices}
S.~Newman,
\textit{Building Microservices}. O'Reilly Media, 2015.

\bibitem{barroso2013datacenter}
L.~A.~Barroso, J.~Clidaras, and U.~H\"{o}lzle,
``The Datacenter as a Computer,''
\textit{Synthesis Lectures on Computer Architecture}, vol.~8, no.~3, 2013.

\bibitem{zhang2020intelligent}
P.~Zhang, M.~Zhou, and X.~Wang,
``An Intelligent Optimization Method for Optimal VM Allocation in Cloud Data Centers,''
\textit{IEEE Trans. Autom. Sci. Eng.}, vol.~17, no.~4, pp.~1725--1735, 2020.

\bibitem{mirjalili2016woa}
S.~Mirjalili and A.~Lewis,
``The Whale Optimisation Algorithm,''
\textit{Advances in Engineering Software}, vol.~95, pp.~51--67, 2016.

\bibitem{storn1997de}
R.~Storn and K.~Price,
``Differential Evolution --- A Simple and Efficient Heuristic for Global
Optimization over Continuous Spaces,''
\textit{Journal of Global Optimization}, vol.~11, no.~4, pp.~341--359, 1997.

\bibitem{lai2025energy}
X.~Lai, X.~Zheng, and W.~Lu,
``Energy-aware VM Migration based on Dual Thresholds,''
in \textit{Proc. IEEE ICBDSE}, 2025.

\bibitem{coffman1984bin}
E.~G.~Coffman~Jr., M.~R.~Garey, and D.~S.~Johnson,
``Approximation Algorithms for Bin Packing,''
in \textit{Approximation Algorithms for NP-Hard Problems}, PWS Publishing, 1996.

\bibitem{nguyen2022loadbalancing}
Q.-M.~Nguyen, L.-A.~Phan, and T.~Kim,
``Load-Balancing of Kubernetes-Based Edge Computing Infrastructure,''
\textit{Sensors}, vol.~22, no.~2869, 2022.

\bibitem{saxena2022secure}
D.~Saxena et al.,
``A Secure and Multiobjective VM Placement Framework for Cloud Data Center,''
\textit{IEEE Systems Journal}, vol.~16, no.~2, pp.~3163--3174, 2022.

\end{thebibliography}

\end{document}